%!TEX root = ../main.tex
\chapter{Introduction to R for Statistics}
\label{ch:r_intro}

\begin{intuition}
If Python is a general-purpose Swiss army knife, \textbf{R}\index{R (language)} is a
precision instrument designed specifically for statistics. Created by
statisticians for statisticians, R excels at data analysis, hypothesis
testing, and visualization. Many researchers in biology, ecology, and
social sciences use it daily.
\end{intuition}

% ----------------------------------------------------------------
\section{Why R?}

\begin{definition}[The R language]
\textbf{R} is a free and open-source programming language specialized in statistical
computing and graphical representation. It has an ecosystem of more than
20,000 packages on the CRAN\index{CRAN} repository.
\end{definition}

\begin{center}
\begin{tikzpicture}[
  langbox/.style={draw, rounded corners, minimum width=3.5cm, minimum height=2.8cm,
                  align=center, font=\small},
  feat/.style={font=\footnotesize, align=left},
]
  \node[langbox, fill=blue!60, text=white] (python) at (0,0)
    {\textbf{\Large Python}};
  \node[feat, below=0.2cm of python.south, anchor=north] {
    $\bullet$ General-purpose\\
    $\bullet$ Machine learning\\
    $\bullet$ Automation\\
    $\bullet$ NumPy, Pandas, SciPy
  };

  \node[langbox, fill=green!50!black, text=white] (r) at (7,0)
    {\textbf{\Large R}};
  \node[feat, below=0.2cm of r.south, anchor=north] {
    $\bullet$ Native statistics\\
    $\bullet$ Hypothesis testing\\
    $\bullet$ Visualization (ggplot2)\\
    $\bullet$ Bioconductor, tidyverse
  };

  \draw[<->, thick, >=stealth] (python.east) -- (r.west)
    node[midway, above, font=\bfseries] {Complementary};
\end{tikzpicture}
\end{center}

% ----------------------------------------------------------------
\section{Variables and vectors}
\index{vector (R)}\index{c() (R)}

In R, the usual assignment operator is \ri{<-} (arrow). The basic structure
is the \textbf{vector}, created with the \ri{c()} function.

\begin{codeblock}[title=\textbf{R}]
\begin{minted}{r}
# Assignment
x <- 42
name <- "physics"

# Numeric vector
measurements <- c(12.3, 14.1, 11.8, 13.5, 12.9)
print(measurements)

# Vectorized operations
measurements_cm <- measurements * 100
print(measurements_cm)
\end{minted}
\end{codeblock}

\begin{codeoutput}
\begin{verbatim}
[1] 12.3 14.1 11.8 13.5 12.9
[1] 1230 1410 1180 1350 1290
\end{verbatim}
\end{codeoutput}

\begin{warning}
In R, indexing starts at \textbf{1} (not 0 as in Python).
Thus, \ri{measurements[1]} returns the \emph{first} element, i.e.\ 12.3.
\end{warning}

\begin{codeblock}[title=\textbf{R}]
\begin{minted}{r}
# Indexing (starts at 1!)
print(measurements[1])    # First element
print(measurements[2:4])  # Elements 2 to 4 (inclusive!)

# Sequences and repetitions
indices <- 1:10
print(indices)
repetitions <- rep(0, 5)
print(repetitions)
\end{minted}
\end{codeblock}

\begin{codeoutput}
\begin{verbatim}
[1] 12.3
[1] 14.1 11.8 13.5
 [1]  1  2  3  4  5  6  7  8  9 10
[1] 0 0 0 0 0
\end{verbatim}
\end{codeoutput}

% ----------------------------------------------------------------
\section{Data frames in R}
\index{data.frame (R)}

\begin{codeblock}[title=\textbf{R}]
\begin{minted}{r}
# Create a data frame
experiment <- data.frame(
  subject = c("A", "B", "C", "D", "E"),
  dose    = c(10, 20, 30, 40, 50),
  effect  = c(2.1, 4.5, 5.8, 7.2, 8.9)
)
print(experiment)
str(experiment)
\end{minted}
\end{codeblock}

\begin{codeoutput}
\begin{verbatim}
  subject dose effect
1       A   10    2.1
2       B   20    4.5
3       C   30    5.8
4       D   40    7.2
5       E   50    8.9
'data.frame':	5 obs. of  3 variables:
 $ subject: chr  "A" "B" "C" "D" ...
 $ dose   : num  10 20 30 40 50
 $ effect : num  2.1 4.5 5.8 7.2 8.9
\end{verbatim}
\end{codeoutput}

% ----------------------------------------------------------------
\section{Basic statistical functions}
\index{mean (R)}\index{sd (R)}\index{var (R)}\index{median (R)}\index{summary (R)}

R natively includes all essential statistical functions.

\begin{codeblock}[title=\textbf{R}]
\begin{minted}{r}
grades <- c(12, 15, 8, 14, 16, 11, 13, 9, 17, 10)

cat("Mean        :", mean(grades), "\n")
cat("Median      :", median(grades), "\n")
cat("Std. dev.   :", sd(grades), "\n")
cat("Variance    :", var(grades), "\n")
cat("Min / Max   :", min(grades), "/", max(grades), "\n")

# Full summary
summary(grades)
\end{minted}
\end{codeblock}

\begin{codeoutput}
\begin{verbatim}
Mean        : 12.5
Median      : 12.5
Std. dev.   : 3.027650
Variance    : 9.166667
Min / Max   : 8 / 17
   Min. 1st Qu.  Median    Mean 3rd Qu.    Max.
   8.00   10.25   12.50   12.50   14.75   17.00
\end{verbatim}
\end{codeoutput}

% ----------------------------------------------------------------
\section{Statistical tests}
\index{t.test (R)}\index{cor.test (R)}

\begin{definition}[Student's t-test]
The \textbf{Student's t-test}\index{t-test} compares the means of two groups
to determine whether the observed difference is statistically significant
(p-value $< 0.05$).
\end{definition}

\begin{codeblock}[title=\textbf{R}]
\begin{minted}{r}
# Two groups: treatment vs control
treatment <- c(23.1, 25.4, 22.8, 26.1, 24.5, 27.0)
control   <- c(19.2, 20.1, 18.5, 21.3, 19.8, 20.5)

result <- t.test(treatment, control)
print(result)
\end{minted}
\end{codeblock}

\begin{codeoutput}
\begin{verbatim}
	Welch Two Sample t-test

data:  treatment and control
t = 5.1893, df = 9.5182, p-value = 0.0004507
alternative hypothesis: true difference in means is not equal to 0
95 percent confidence interval:
 2.633574 6.566426
sample estimates:
mean of x mean of y
 24.81667  20.21667
\end{verbatim}
\end{codeoutput}

\begin{codeblock}[title=\textbf{R}]
\begin{minted}{r}
# Correlation test
x <- c(1, 2, 3, 4, 5, 6, 7, 8)
y <- c(2.1, 4.3, 5.8, 8.2, 9.7, 12.1, 13.8, 16.2)
cor.test(x, y)
\end{minted}
\end{codeblock}

\begin{codeoutput}
\begin{verbatim}
	Pearson's product-moment correlation

data:  x and y
t = 42.085, df = 6, p-value = 4.386e-08
95 percent confidence interval:
 0.9971346 0.9999227
sample estimates:
      cor
0.9984112
\end{verbatim}
\end{codeoutput}

% ----------------------------------------------------------------
\section{Linear regression with \ri{lm()}}
\index{lm() (R)}\index{linear regression}

\begin{codeblock}[title=\textbf{R}]
\begin{minted}{r}
# Linear regression: effect as a function of dose
model <- lm(effect ~ dose, data = experiment)
summary(model)
\end{minted}
\end{codeblock}

\begin{codeoutput}
\begin{verbatim}
Coefficients:
            Estimate Std. Error t value Pr(>|t|)
(Intercept)  0.14000    0.34351   0.408  0.71062
dose         0.17200    0.01033  16.649  0.00048 ***
---
Residual standard error: 0.3742 on 3 degrees of freedom
Multiple R-squared:  0.9893,	Adjusted R-squared:  0.9857
\end{verbatim}
\end{codeoutput}

The coefficient $R^2 = 0.989$ indicates an excellent linear fit:
the effect increases by 0.172 units per unit of dose.

% ----------------------------------------------------------------
\section{Basic plots in R}
\index{plot (R)}\index{hist (R)}\index{boxplot (R)}

\begin{codeblock}[title=\textbf{R}]
\begin{minted}{r}
# Scatter plot with regression line
plot(experiment$dose, experiment$effect,
     xlab = "Dose (mg)", ylab = "Effect",
     main = "Dose-effect relationship", pch = 19, col = "blue")
abline(model, col = "red", lwd = 2)

# Histogram
hist(rnorm(1000), breaks = 30, col = "lightblue",
     main = "Simulated normal distribution",
     xlab = "Value")

# Box plot
boxplot(treatment, control,
        names = c("Treatment", "Control"),
        col = c("coral", "lightgreen"),
        main = "Group comparison")
\end{minted}
\end{codeblock}

% ----------------------------------------------------------------
\section{Python vs R comparison}

\begin{example}[Same analysis in both languages]

\begin{codeblock}[title=\textbf{Python}]
\begin{minted}{python}
import numpy as np
from scipy import stats

data = [12, 15, 8, 14, 16, 11, 13, 9, 17, 10]
print(f"Mean: {np.mean(data):.2f}")
print(f"Std. dev.: {np.std(data, ddof=1):.2f}")

# t-test
group_a = [23.1, 25.4, 22.8, 26.1, 24.5, 27.0]
group_b = [19.2, 20.1, 18.5, 21.3, 19.8, 20.5]
t_stat, p_val = stats.ttest_ind(group_a, group_b)
print(f"p-value: {p_val:.6f}")
\end{minted}
\end{codeblock}

\begin{codeblock}[title=\textbf{R}]
\begin{minted}{r}
data <- c(12, 15, 8, 14, 16, 11, 13, 9, 17, 10)
cat("Mean:", mean(data), "\n")
cat("Std. dev.:", sd(data), "\n")

# t-test
group_a <- c(23.1, 25.4, 22.8, 26.1, 24.5, 27.0)
group_b <- c(19.2, 20.1, 18.5, 21.3, 19.8, 20.5)
result <- t.test(group_a, group_b)
cat("p-value:", result$p.value, "\n")
\end{minted}
\end{codeblock}
\end{example}

\begin{remark}
Note that \ri{sd()} in R computes the \emph{corrected} standard deviation ($n-1$) by default,
like \py{np.std(data, ddof=1)} in Python. The function \py{np.std()} without
\py{ddof=1} divides by $n$, which gives a different result.
\end{remark}

% ----------------------------------------------------------------
\begin{errmark}{Common errors in R}
\begin{enumerate}
  \item \textbf{Indexing from 1}\index{R indexing} --- In R, \ri{x[1]} is the
    \emph{first} element. In Python, \py{x[1]} is the \emph{second}. This difference
    is the source of many errors when switching between languages.
  \item \textbf{\ri{<-} vs \ri{=}} --- In R, \ri{<-} is preferred for assignment.
    The \ri{=} sign also works in most cases, but \ri{<-} is the
    standard convention. Be careful not to write \ri{x < -5} (comparison!) instead
    of \ri{x <- 5} (assignment).
  \item \textbf{Unexpected factors} --- R sometimes automatically converts strings
    to \texttt{factor}. Use \ri{stringsAsFactors = FALSE} in
    \ri{data.frame()} if needed.
  \item \textbf{Forgetting \ri{na.rm = TRUE}} --- Functions like \ri{mean()} return
    \ri{NA} if the vector contains missing values. Add \ri{na.rm = TRUE}.
\end{enumerate}
\end{errmark}

% ----------------------------------------------------------------
\begin{projet}{Comparative statistical analysis in R}
Carry out a complete study in R:
\begin{enumerate}
  \item Create a data frame with two experimental groups (15 measurements each),
    simulated with \ri{rnorm(15, mean=..., sd=...)}.
  \item Compute the descriptive statistics for each group (\ri{mean}, \ri{sd},
    \ri{median}).
  \item Perform a Student's t-test to compare the two groups.
  \item Draw side-by-side box plots.
  \item Add a third variable and perform a linear regression with \ri{lm()}.
  \item Interpret the results: is the difference significant?
    Is the linear model appropriate ($R^2$)?
\end{enumerate}
\end{projet}

% ----------------------------------------------------------------
\section{Exercises}

\begin{exercise}[$\star$]
Create a vector containing the average monthly temperatures of a city
(12 values). Compute the mean, median, and standard deviation. Which month is the
hottest? Use \ri{which.max()}.
\end{exercise}

\begin{exercise}[$\star$]
Create a data frame with the columns \ri{element}, \ri{symbol}, and
\ri{atomic\_mass} for 5 chemical elements. Display the summary with
\ri{summary()} and access the \ri{atomic\_mass} column with \ri{\$}.
\end{exercise}

\begin{exercise}[$\star\star$]
Generate 100 random values following a normal distribution ($\mu = 50$, $\sigma = 10$)
with \ri{rnorm()}. Plot a histogram and overlay the theoretical curve with
\ri{curve(dnorm(x, 50, 10), add = TRUE)}.
\end{exercise}

\begin{exercise}[$\star\star$]
Simulate an experiment: generate two samples of size 30, one with mean
100 and the other with mean 105 (same standard deviation 15). Perform a t-test.
Repeat 1000 times and count the percentage of times $p < 0.05$
(this is the \textbf{statistical power}\index{statistical power} of the test).
\end{exercise}

\begin{exercise}[$\star\star\star$]
Load the built-in dataset \ri{iris} in R. Carry out a complete analysis:
descriptive statistics by species, ANOVA test with \ri{aov()} to compare
sepal lengths between species, and linear regression between petal length and
petal width. Produce 3 relevant plots.
\end{exercise}

% ----------------------------------------------------------------
\begin{keyformulas}{Chapter 9 --- R: the essentials}
\begin{tabular}{@{}p{5cm} p{9cm}@{}}
\toprule
\textbf{Concept} & \textbf{R syntax} \\
\midrule
Assignment & \ri{x <- value} \\
Vector & \ri{c(1, 2, 3)} \\
Sequence & \ri{1:10} or \ri{seq(0, 1, by=0.1)} \\
Data frame & \ri{data.frame(col1=..., col2=...)} \\
Mean & \ri{mean(x)} \\
Standard deviation & \ri{sd(x)} \\
Summary & \ri{summary(x)} \\
t-test & \ri{t.test(group1, group2)} \\
Correlation & \ri{cor.test(x, y)} \\
Regression & \ri{lm(y \~{} x, data=df)} \\
Plot & \ri{plot(x, y)}, \ri{hist(x)}, \ri{boxplot(x)} \\
\bottomrule
\end{tabular}
\end{keyformulas}
